% Generated draft: requires mathematical and source review.
\documentclass[11pt]{article}
\usepackage[margin=25mm]{geometry}
\usepackage{fontspec}
\setmainfont{Latin Modern Roman}
\usepackage{amsmath,amssymb,mathtools,unicode-math}
\setmathfont{Latin Modern Math}
\usepackage{enumitem,needspace,xurl,hyperref}
\setlength{\parindent}{0pt}
\setlength{\parskip}{5pt}
\setlength{\emergencystretch}{3em}
\newcommand{\N}{\mathbb{N}}
\newcommand{\Z}{\mathbb{Z}}
\newcommand{\Q}{\mathbb{Q}}
\newcommand{\R}{\mathbb{R}}
\newcommand{\C}{\mathbb{C}}
\newcommand{\F}{\mathbb{F}}
\newcommand{\A}{\mathbb{A}}
\newcommand{\PP}{\mathbb{P}}
\newcommand{\E}{\mathbb{E}}
\newcommand{\eps}{\varepsilon}
\newcommand{\dd}{\,\mathrm{d}}
\begin{document}
\section*{UnsolvedMath ID: 30000398}
\textbf{Source problem code:} catalogue entry
\subsection{Definitions and mathematical statement}
The historical question is whether the balance equation of canonical momentum in continuum mechanics can be formulated without specifying an energy density. The supplied source-verification note explicitly reports an affirmative answer in the official report, with the essential qualification that such a formulation is not formally independent of the usual balance equations. This record is therefore retained as an answered question with that attributed qualification, not promoted to a currently open conjecture.

To distinguish the requested objects abstractly, write \(\mathcal P\) for an intended canonical momentum density, \(\mathcal B\) for its flux and \(\mathcal F\) for a source term. A local balance could schematically be written
\[
 \partial_t\mathcal P+\operatorname{Div}\mathcal B=\mathcal F.
\]
The question concerns whether the intended balance can be formulated without first choosing a constitutive energy density \(W\). This display is only a general notation for a balance, not the actual canonical-momentum formula from the report. Neither the continuum fields, definition of canonical momentum, flux, forces, material or spatial coordinate convention, nor the energy variables are supplied. In particular canonical momentum is not silently identified with ordinary mass-times-velocity momentum or with a chosen Lagrangian derivative.

The reported affirmative answer must be read together with the dependence qualification: formulating the balance without specifying \(W\) does not establish an extra independent equation beyond the usual balances. It also does not assert that \(\mathcal P\) is conserved without sources, that arbitrary constitutive choices are equivalent, or that the statement provides a closed system determining the motion. The precise algebraic or differential dependence on the usual balance laws is absent from the record and cannot be reconstructed.

The literature triage checked 2026-08-21 instead calls the item an explicit unresolved question and reports no located authoritative general resolution, while the later verification reports that the report itself answers it. The registry status remains open, the statement-status is unrecoverable, and the cleaned field is null. These conflicting curation labels are recorded as metadata; the explicit answered-status qualification is preserved without claiming an independently checked derivation.
\subsection{Short English statement}
The report is said to answer yes: canonical-momentum balance can be formulated without choosing an energy density, but the formulation is not formally independent of the usual balances. The actual formula and its dependence are absent. The draft preserves this answered historical question and flags the contradictory open-status metadata.
\subsection{Sources}
\begin{itemize}
\item Local UnsolvedMath registry, permanent ID 30000398, source code catalogue entry, ``Answered Canonical-Momentum Balance Question,'' including the original question, null cleaned statement and verification note reporting the affirmative answer and dependence qualification.
\item \emph{Mechanics of Materials (2006)}, Oberwolfach Reports, pp.\ 187--252, DOI 10.4171/owr/2006/04: \url{https://doi.org/10.4171/owr/2006/04}. This is the sole supplied literature source, labelled ``original Oberwolfach report.''
\item Literature triage checked 2026-08-21; local record updated 2026-09-16. The unresolved assessment conflicts with the explicit answered-status verification. No external source has been visited or independently checked here.
\end{itemize}
\end{document}
